\chapter{Cwiczenia i propozycje samodzielnej pracy}

\section{Dlaczego warto robic cwiczenia}
Sama lektura dokumentacji daje podstawy, ale prawdziwe zrozumienie przychodzi dopiero,
gdy samodzielnie zmieniasz kod i obserwujesz efekt. Dlatego na koniec warto przejsc
przez kilka malych zadan.

\section{Cwiczenie 1: prosta strona wizytowka}
Cel: utworz \mc{chunk} generujacy pojedyncza strone glowna z naglowkiem, opisem i linkiem.

Spróbuj samodzielnie:

\begin{itemize}
  \item dodac \mc{title} i \mc{metaDescription},
  \item zdefiniowac trzy zmienne tekstowe,
  \item wstawic je do \mc{content}.
\end{itemize}

\section{Cwiczenie 2: lista kart}
Cel: zbuduj tablice trzech obiektow z polami \mc{title}, \mc{description}, \mc{button}
i wygeneruj z niej karty w petli \mc{for}.

To cwiczenie utrwala:

\begin{itemize}
  \item tablice,
  \item obiekty,
  \item indeksowanie,
  \item dynamiczne \mc{content} w petli.
\end{itemize}

\section{Cwiczenie 3: filtr opublikowanych wpisow}
Cel: przygotuj tablice wpisow z polem \mc{published} i pokaz tylko te, ktore sa gotowe do
publikacji. Dodaj tez tekst ``brak wpisow'', jesli zadny nie przejdzie filtra.

\section{Cwiczenie 4: funkcja etykietujaca}
Cel: napisz funkcje, ktora na podstawie liczby wyswietlen zwraca etykiete,
np. \mc{"Nowy"}, \mc{"Popularny"}, \mc{"Hit"}. Uzyj tej funkcji w HTML-u.

\section{Cwiczenie 5: wczytanie JSON}
Cel: przenies dane wpisow z kodu MoonChunk do pliku JSON i wczytaj je przez \mc{data(...)}.
Sprawdz przez \mc{print(...)}, czy dane wygladaja tak, jak oczekujesz.

\section{Cwiczenie 6: wspolne metadata}
Cel: wydziel chunk \mc{Metadata}, zaimportuj go do pliku glownego i nadpisz \mc{title}
dla jednej konkretnej strony.

\section{Cwiczenie 7: rekurencja}
Cel: napisz funkcje rekurencyjna, ktora wypisuje liczby od \mc{0} do \mc{5}. Zadbaj o
czytelny warunek stopu.

\section{Cwiczenie 8: jawny blok zasiegu}
Cel: zadeklaruj zmienna \mc{x} na zewnatrz, a potem utworz jawny blok \mc{\{ ... \}},
w ktorym pojawi sie kolejne \mc{x}. Sprawdz przez \mc{print(...)}, ze obie wartosci sa
od siebie niezalezne.

\section{Jak pracowac z cwiczeniami}
Najlepiej wykonywac je w kolejnosci. Po kazdym z nich odpowiedz sobie na trzy pytania:

\begin{enumerate}
  \item co juz rozumiem lepiej niz przed chwila,
  \item gdzie musialem uzyc \mc{print(...)},
  \item ktory fragment nadal budzi watpliwosci.
\end{enumerate}

To pomaga nie tylko rozwiazywac zadania, ale tez świadomie budowac intuicje jezykowa.
